\chapter{Angles and Trigonometric Functions}\label{ch:rational-circle-trigonometry}

The preceding chapter constructed circles, arc lengths, sector areas, and
the sphere--cylinder comparisons without trigonometric functions.  We now
use sector area to assign an angle coordinate to a point of the circle.
The new objects are functions of that coordinate:
\[
 S(x)=\sin(\pi x),\qquad C(x)=\cos(\pi x).
\]
These names refer to the computations below, not to functions imported
with an existing theory of differentiation.

\section{The sector-area coordinate}
\label{sec:sector-angle-coordinate}
Recall the rational circle chart and its sector-area computation from
Chapter~\ref{ch:circle-sphere}:
\[
 P(t)=\left(\frac{1-t^2}{1+t^2},\frac{2t}{1+t^2}\right),
 \qquad A(1)=\pi/4.
\]
A full turn has normalized parameter length $2$ and a quarter turn has
length $1/2$.  For rational $0\le x\le1/2$, the point with sector area
$\pi x/2$ is to have coordinates $(C(x),S(x))$.

The geometric estimates already proved are
\[
 \frac{v-u}{1+v^2}\le A(v)-A(u)\le\frac{v-u}{1+u^2},
 \qquad \frac{v-u}{2}\le A(v)-A(u)\le v-u
\]
for rational $0\le u<v\le1$.  The chart therefore has interval
continuity and a linear inverse-separation bound.  The target
$\pi x/2=2xA(1)$ is enclosed by its endpoint values.  The trisection
construction of Chapter~\ref{ch:foundations} computes a unique parameter
$t$ such that
\[
 A(t)=\pi x/2.
\]
Evaluation of $P$ at this computed input gives $C(x),S(x)$.  At $x=0$
and $x=1/2$ use $t=0,1$ directly.  Reflections in the coordinate axes and
whole turns give the computation at every rational input.

The conventional name $A(t)=\arctan t$ records the half-angle nature of
this rational chart: the arc length from $(1,0)$ to $P(t)$ is $2A(t)$.
No integral definition of arctangent is needed.  The rational rectangle
bounds for $1/(1+t^2)$ that occur in the sector-area construction will
later give an equivalent integral computation.

\section{Dyadic rotations as another computation}
\label{thm:finite-quarter-turn}
If $(c,s)$ is a unit-circle point in the first quadrant, its half-angle
point is
\[
 (a,b)=\left(\sqrt{(1+c)/2},\sqrt{(1-c)/2}\right).
\]
Finite algebra gives $a^2+b^2=1$, $a^2-b^2=c$, and $2ab=s$; the
nonnegative sign in the last identity follows from $s\ge0$.  Thus the
coordinate rotation matrix of $(a,b)$ squares to that of $(c,s)$.
Additivity and invariance of sector area show that its angle coordinate
is half that of the parent point.

Repeated halving of the quarter turn, followed by finite rotation products,
therefore constructs all dyadic angle samples by nested square roots.
They agree with the inverse-sector construction by its separation and
uniqueness.  This agreement is a geometric theorem between two computations,
not an appeal to an already completed trigonometric function.

\section{Addition laws and continuity}
Multiplying the coordinate rotation matrices and using the additivity of
sector area gives
\begin{align*}
 S(x+y)&=S(x)C(y)+C(x)S(y),\\
 C(x+y)&=C(x)C(y)-S(x)S(y),\\
 C(x)^2+S(x)^2&=1.
\end{align*}
The endpoint values are
$S(0)=0$, $C(0)=1$, $S(1/2)=1$, and $C(1/2)=0$.
Reflections give $S(-x)=-S(x)$ and $C(-x)=C(x)$; both have period $2$.

A coordinate difference is bounded by chord length, and chord length by
arc length.  The arc-length computation from Chapter~\ref{ch:circle-sphere}
and the bound $\pi<4$ yield
\[
 |S(x)-S(y)|\le4|x-y|,\qquad
 |C(x)-C(y)|\le4|x-y|.
\]
These hold for every rational pair, splitting across quadrant boundaries
when necessary.  Thus dyadic approximation is another terminating way to
evaluate any rational angle.  They also give evaluation at computed inputs
on specified bounded intervals.

\section{Small angles before differentiation}
\label{sec:geometric-small-angle}
For $0<h\le1/4$, chord, sector, and tangent comparisons give
\[
 S(h)\le\pi h\le S(h)/C(h).
\]
For the left inequality, the vertical displacement is at most the chord,
which is at most the arc length.  For the right inequality, twice the
sector area is at most twice the area of the enclosing tangent triangle.
The geometric point lies between the positive horizontal axis and its
bisector, so $C(h)>0$ with a positive computable lower bound.

The half-angle identity and chord bound give
\[
 0\le1-C(h)=2S(h/2)^2\le8h^2.
\]
Multiplying the upper tangent comparison by $C(h)$ and using $\pi<4$
yields
\[
 0\le\pi h-S(h)\le\pi h(1-C(h))\le32h^3.
\]
Odd and even symmetry therefore give, for $|h|\le1/4$,
\[
 |S(h)-\pi h|\le32|h|^3,\qquad |1-C(h)|\le8h^2.
\]
The second estimate also holds for all rational $h$, by the global chord
bound and half-angle identity.  These finite geometric errors will prove
both that the convex derivative computations terminate and that their
values are $\pi C$ and $-\pi S$.

\section{Input to integration}
One of the first areas computed in two ways will be
\[
 \int_0^{1/2}S(x)\,dx.
\]
Equal-angle samples use the nested-radical dyadic construction.  Unequal
samples use the rational circle parameter $t$ and the sector coordinate
$2A(t)/\pi$.  Their comparison uses the explicit continuity bounds and
finite refinements.  The later derivative computation and FTC identify
the answer as $1/\pi$.

The distinction between the chapters is now explicit: the preceding one
computes particular geometric magnitudes, while this one constructs angle
functions and compares their descriptions.  General integration organizes
both sorts of computations only after they have been established.
